%! program = pdflatexmk

% VoteSecure FAQ
% Copyright (C) 2025-26 Free & Fair

\documentclass[12pt,letter]{article}

%%% PACKAGES
\usepackage[in]{fullpage}
\usepackage{tabularray}
\usepackage[pdftex,bookmarks,colorlinks,breaklinks]{hyperref}
\usepackage{amsmath}
\usepackage{amssymb}
\usepackage{listings}
\usepackage{verbatim}
\usepackage{lmodern}
\usepackage[T1]{fontenc}
\setlength {\marginparwidth}{2cm}
\usepackage{todonotes}
\usepackage{enumitem}
\hypersetup{linkcolor=blue,citecolor=blue,filecolor=black,urlcolor=blue}
\parindent 0pt
\textheight 9in
\topmargin 0in
\setlength{\parskip}{10pt plus 1pt minus 1pt}

\usepackage{tocloft}
\setlength\cftparskip{2pt}

% format FAQ date in British English, but without the ordinal suffix
\usepackage[en-GB]{datetime2}
\DTMlangsetup[en-GB]{ord=omit}

\renewcommand{\subsectionautorefname}{section}
\renewcommand{\subsubsectionautorefname}{section}

% by default, don't put headings on tables
\DefTblrTemplate{firsthead, middlehead,lasthead}{default}{}

\pagestyle{empty}

% hyphenation
\hyphenation{Vote-Secure}

%%% BEGIN DOCUMENT
\begin{document}

\begin{center}
{\Large \textbf{VoteSecure FAQ}}
\end{center}
\vspace{-6pt}

\textit{Free \& Fair} \hspace{\fill} \textit{\DTMtoday}

\vspace{-12pt}

\rule{\textwidth}{1pt}

The concept of mobile voting naturally raises many questions, both about the idea itself and about the role of the VoteSecure library in the implementation of any publicly-available mobile voting system. This document contains our answers to questions we have been asked, and will be updated with future questions and answers as they arise.

% Q & A go here, we can format them in a consistent manner later on

\section{General Project Questions}

\begin{itemize}

    \item \emph{Q:} Where should I start reading?

    \emph{A:} It's best to start at the project page on our website and go from there. \\
    \href{https://freeandfair.us/votesecure/}{https://freeandfair.us/votesecure/}

    \item \emph{Q:} Is there more information coming out soon?

    \emph{A:} Every time we merge new content into the \texttt{main} branch of the project it is mirrored to GitHub.  So, while we have a few active branches of work (e.g., on mechanizing the protocol in Tamarin) that will not be public until they are merged into \texttt{main}, the vast majority of the work is public already, and it will all be public when it is complete.

    \item \emph{Q:} Is most or all of the technology developed by the Danish company Assembly Voting gone? Are you making something different?

    \emph{A:} Yes, all of Assembly Voting's work is gone.  We analyzed their cryptographic protocol and implementation; \href{https://github.com/FreeAndFair/Transparency/blob/master/Project_Artifacts/Tusk_2024_Review/cryptographic-protocol-review-and-gap-analysis.pdf}{our report summarizing that analysis} is publicly available in our \href{https://github.com/FreeAndFair/Transparency}{Transparency repository} on GitHub.

    When this project started we had originally thought that portions of that protocol might be reusable, but as we performed a more detailed security analysis of the protocol we found so many flaws that it became clear we had to start from scratch.

    \item \emph{Q:} Will white hat hackers be given legal protection and a reasonable lead time prior to the system being used in any election?  If so, will there be any publication, NDA, or other constraints on what the hackers may say/write?

    \emph{A:} We are only developing the core cryptographic library, and we cannot speak for system implementations; however, we welcome security analysis of both the cryptographic protocol and (when available) its implementation. The project repository contains \href{https://github.com/FreeAndFair/VoteSecure/blob/main/SECURITY.md}{instructions for reporting security issues} that follow standard \href{https://en.wikipedia.org/wiki/Coordinated_vulnerability_disclosure}{coordinated vulnerability disclosure (CVD)} guidelines. Other than requesting that they follow CVD guidelines, we will not place any restrictions on what anybody can write or say.

\end{itemize}


\section{Cryptography}

\begin{itemize}

    \item \emph{Q:} Are you using a mixnet?

    \emph{A:} We are using a mixnet in the first version, because that's the most straightforward way to produce CVRs that can be used to generate normal paper ballots. The election officials (EOs) Tusk has been working extensively with want the back end of the system to either be a human-centric remarking workflow exactly like they use today, or to produce PDFs of normal paper ballots that are readable by their existing scanners.  Note that we are only implementing the cryptographic protocol, not this back end, in our current work.

    \item \emph{Q:} Do you support homomorphic tallying?

    \emph{A:} While a homomorphic tally is possible for our protocol, and is a variant we have examined closely, we are not implementing it at this time.  Such a back end has an entirely different workflow for EOs---one that they are not familiar or comfortable with as of now, as we understand it. It also presents issues with respect to certain ballot features, such as write-in candidates.

    \item \emph{Q:} Do you use a Benaloh-like challenge as a part of the voter's user experience as seen in the original AV protocol?

    \emph{A:} No, we do not use a Benaloh-like challenge; in our protocol the voter checks the decrypted contents of the actual encrypted ballot they submitted before deciding whether to cast it, rather than challenging it and then always submitting a new encrypted ballot with the same contents as in the AV protocol. As you will see in our analysis, the AV Benaloh-like challenge had fundamental security flaws, which perhaps contributed to the positive outcomes in their usability studies. We have had to think hard about the voter experience in order to enable a user experience that fulfills E2E-V requirements and yet is usable by and understandable to all voters.

    \item \emph{Q:} Is the verification by voters of ``cast as intended'' voluntary, or mandatory?

    \item \emph{A:} The current version of the protocol does not require voters to perform the verification step. It is possible to make the verification step mandatory at the protocol level, or to configure a particular jurisdiction's voting application such that the verification step is ``effectively mandatory'' (i.e., the voting application will not proceed to the casting step without the verification step having occurred at least once). However, we do not believe that making the verification step mandatory at the protocol level provides any additional guarantees over making it effectively mandatory at the voting application level; there is no way for the protocol to guarantee that a voter actually performs the verification step properly (as opposed to just taking the required actions and ignoring the actual verification information presented). We presume that any voter who would perform the check properly if it were mandatory at the protocol level would do the same if it were effectively mandatory due to voter application flow.

\end{itemize}

\section{Election Officials}

\begin{itemize}

    \item \emph{Q:} Aren't there real concerns about the security (the ordinary kind) of the IT systems to facilitate elections, and the feasibility of ordinary EOs doing the key management and key usage?

    \emph{A:} While key management is a necessary part of most systems that perform cryptography, we hold low concern about EOs ability to generate, manage, and use keys in our system.  The protocol itself identifies key management problems and threats as a part of its overall CONOPS and threat model, as mistakes and adversaries focused on key material are important. However, given modern key generation and management systems that are now widely deployed in public, including key managers in OSs, browsers, and password managers, as well as more esoteric, secure, and somewhat less user-friendly devices like Yubikeys and Google Titan keys, we are in a much better place today than we were just a few years ago.

    \item \emph{Q:} I liked (in theory) the basic idea of a digital process very similar to the paper absentee ballot process: the ballot is cloaked and accompanied by an affidavit; and it is the EO's responsibility to decide (based on the affidavit and other stuff) whether each ballot is admissible, and only if so is the cloaked ballot queued up for counting. I don't know whether the digital version of that approach (including the same literal affidavit document as in paper AB) can be done in the scheme you're using, but I am curious if it is a goal, and if so whether it can be achieved using SW that is usable by ordinary EOs, and deployed on IT systems that can be protected by EO IT staff and systems.

    \emph{A:} The goal is indeed to support that kind of workflow as a part of our protocol.  It is what EOs want, and they are the domain experts.  Thus, in our protocol the authentication of voters is delegated to a black box component with a generic interface, set of behaviors, and risk characterization.

    \item \emph{Q:} How would a credible recount be conducted?

    \emph{A:} VoteSecure does not tally ballots; it generates paper ballots based on the ballot cryptograms cast by voters, and those ballots are tallied in the same way as other paper ballots in the jurisdiction. The various cryptographic proofs generated by VoteSecure during ballot cryptogram generation, shuffling, and decryption ensure that there is only one possible set of generated paper ballots for any set of ballot cryptograms, so there is nothing to ``recount'' within the protocol. Any recount would therefore be a recount of the generated paper ballots subject to the policies and procedures used for such recounts, including issues like chain of custody for the generated ballots.

    \item \emph{Q:} What are your thoughts on threat modeling the potential substantial harm to legitimacy of elections if there are notable attacks that are detected and publicized, even with E2E-V?

    \emph{A:} We believe that such modeling is important, and that a good understanding of the potential harms is necessary before widespread deployment of voting systems based on VoteSecure. However, while some of us have studied political science and other social sciences, none of us are experts in those fields. Modeling the possible effects that hypothetical notable attacks on future election systems might have on human behavior and the legitimacy of democratic systems is far outside our area of expertise, and should be done by more qualified parties.

    \item \emph{Q:} How can authentication be accomplished in a trustworthy fashion, since it is a lamentable fact that US Government at all levels (Federal, State, and especially local Election Officials) completely lack the capability to perform even the fallible types of strong authentication that we have in mainline enterprise IT security?

    \emph{A:} Authentication is deliberately a ``black box'' in VoteSecure; that is, we do not specify any particular method of authentication, only that the results have to be provided to the protocol in a specific format. There are multiple authentication providers in the U.S. that are considered strong enough for purposes such as paying taxes and accessing tax records and other government information (\href{https://id.me}{ID.me}), airport screening and security at other venues (\href{https://clearme.com}{CLEAR}), and more. While these are not as strong as the authentication mechanisms available in countries (e.g., Estonia) where citizens are issued digital credentials in the form of smart cards, we believe that they compare favorably to the authentication mechanisms used in existing remote voting systems (e.g., signature comparison for mail-in ballots).

\end{itemize}

\section{Security}

\begin{itemize}

    \item \emph{Q:} How will the voter be authenticated?

    \emph{A:} Authentication is a ``black box'' in our protocol, meaning that system implementers and EOs must decide upon appropriate authentication processes/providers for their individual implementations and applications. We expect that identity verification systems like \href{https://clearme.com/}{CLEAR} will be used for this purpose, though other approaches (e.g., a digital equivalent to the signed affidavits that accompany mail-in ballots today) can be implemented as well.

    \item \emph{Q:} How will the system defend against distributed denial of service (DDoS) attacks?

    \emph{A:} We are only developing the core cryptographic library, and we cannot speak for system implementations; however, we expect that system implementers will use industry-standard methods for defending Internet-facing systems against DDoS attacks (i.e., \href{https://cloudflare.com/}{Cloudflare} and similar services).

    \item \emph{Q:} How can underfunded and under-resourced local election officials protect their systems from attacks, both external and from insiders? What about ransomware?

    \emph{A:} As our threat model (\href{https://github.com/FreeAndFair/VoteSecure/releases/download/latest/threat-model.pdf}{static PDF}, \href{https://github.com/FreeAndFair/VoteSecure/releases/download/latest/threat-model.html}{interactive HTML}) states, we are explicitly not addressing these sorts of system security issues as part of this project; these will need to be addressed by system implementers. However, we have designed the cryptographic protocol such that (1) voter privacy cannot be compromised without using private key material belonging to a sufficiently large number of election trustees, and (2) ballots cannot be changed in undetectable ways after being submitted by voters.

    \item \emph{Q:} Why is the vulnerability of a consumer-managed, Internet-connected client not reason enough why mobile voting is far too dangerous to be used for public elections?

    \emph{A:} As we note in the answer to the following question, this vulnerability is precisely one of the vulnerabilities mitigated by E2E-V protocols like VoteSecure.

    \item \emph{Q:} Isn't the whole point of E2E-V protocols to mitigate potential cyber-vulnerabilities in client devices and in the servers running elections?

    \emph{A:} While mitigating vulnerabilities in client devices and servers is not the \emph{whole point} of E2E-V protocols (they also mitigate against corrupt election officials,and a number of other threats), it is one of the benefits of an E2E-V protocol. When the protocol is implemented properly, and when voters carry out the ballot check procedure on one or more additional devices, all of a voter's devices would need to be corrupted simultaneously, in a coordinated fashion (i.e., the corrupted devices would need to communicate with each other during the voting and ballot checking processes), to modify that voter's vote.

\end{itemize}

\section{Usability}

\begin{itemize}

    \item \emph{Q:} One issue with E2E-V is that it's not easy for non-technical people to understand why it works. How do you propose to address the trust issue?

    \emph{A:} We have attempted to make the protocol's implementation of end-to-end verifiability straightforward; however, work will clearly need to be done to present it to voters in an approachable, understandable way. This will involve user interface design work for any voting and ballot check applications that implement the protocol, as well as other explanatory resources (documents, presentations, etc.).

    \item \emph{Q:} How can we explain to the average voter, who needs to understand why a remote voting system is secure?

    \emph{A:} See the previous answer (which is specifically about E2E-V, but also applies to security in general).

    \item \emph{Q:} Why not use code voting as a means by which to mitigate adversaries on the voter's client?

    \emph{A:} Code voting mitigates some, but not all, attacks against client devices. In particular, it mitigates attacks related to vote privacy because the device never sees the voter's actual choices. However, code voting introduces considerable infrastructure to handle secure code generation and distribution, which is vulnerable to multiple types of attack. In addition, voters may have more trouble understanding the code voting system and being sure that they have voted the way they intended. All Internet-facing systems involve a balance between security and usability; we believe that VoteSecure enables an appropriate balance, but we realize that others may view that balance differently.

\end{itemize}

\end{document}
